% 【本文件写什么】一级组件 wateros-pseudo-shell（§ 级聚合）
%   - 组件职责与在内核中的位置：内核内嵌伪 shell，用于 bring-up 交互调试。
%   - 聚合 crate os/components/wateros-pseudo-shell/src/lib.rs 的 pub mod 树与对外导出面
%   - 子模块划分、与相邻组件的依赖边界（谁依赖谁、走哪条门面）
%   - 根 feature / 本组件 Cargo.toml feature 如何选用各 impl
%   - 1～2 段综述后，由下方 \input 展开子模块；不在此逐条列举 api trait
% 【事实来源】
%   - os/components/wateros-pseudo-shell/src/lib.rs、Cargo.toml
%   - docs/exports/public-api/wateros-pseudo-shell.md
%   - docs/exports/features/wateros-pseudo-shell.md
%   - os/feature-tree.txt 中 wateros-pseudo-shell 相关段
% 【不写什么】
%   - api-v0 契约细节 → 子目录 api/api-v0.tex
%   - 各 impl 算法/平台细节 → 子目录 impl/*.tex

\section{伪Shell组件}

\code{wateros-pseudo-shell}提供经MMIO UART读行的阻塞式伪shell，用于bring-up阶段验证VFS与用户ELF装载。本节未纳入正文，仅占位；实现细节见\code{os/components/wateros-pseudo-shell/src/lib.rs}。
